\documentclass[11pt,a4paper]{article}
\usepackage[T1]{fontenc}
\usepackage{lmodern,amsmath,amssymb,amsthm,mathtools,microtype}
\usepackage[margin=27mm]{geometry}
\usepackage[colorlinks=true,linkcolor=blue,citecolor=blue,urlcolor=blue]{hyperref}
\newtheorem{theorem}{Theorem}[section]
\newtheorem{proposition}[theorem]{Proposition}
\newtheorem{lemma}[theorem]{Lemma}
\newtheorem{corollary}[theorem]{Corollary}
\newcommand{\Q}{\mathbb Q}\newcommand{\Z}{\mathbb Z}
\DeclareMathOperator{\den}{den}
\DeclareMathOperator{\Tr}{Tr}
\title{Restricted $abc$ hypotheses on the scalar families}
\author{}\date{9 November 2026}
\begin{document}\maketitle
\begin{abstract}
We examine two unconditional restricted estimates suggested by the
attached $abc$ literature survey. For the Frey curves of a sum of even
powers, the reduced denominator of the $j$-invariant is exponentially
large in its logarithmic height. Consequently a polynomial bound for
that denominator cannot hold uniformly in the scalar family. We also
show that the fixed asymmetric-size hypotheses of two subexponential
$abc$ estimates fail along infinitely many required integer witnesses.
These are precise restrictions on applying the cited estimates to the
present construction, not assertions about all possible constructions.
\end{abstract}

\section{The denominator of the Frey invariant}

For $r\in\Q$, let $h(r)$ be its logarithmic projective height and
$\den(r)$ its positive reduced denominator. Fix a positive even integer
$d$. For nonzero coprime integers $A,B$, put
\[
 a=A^d>0,\qquad b=B^d>0,\qquad c=a+b,
 \qquad E_{A,B,d}:y^2=x(x-a)(x+b).
\]
The curve is nonsingular since $abc\ne0$.

\begin{proposition}\label{prop:den}
Let $H=\max(|A|,|B|)$. Then
\[
 j(E_{A,B,d})=\frac{256(a^2+ab+b^2)^3}{a^2b^2c^2}.
\]
The common divisor of this displayed numerator and denominator is
\[
 2^{\min(8,v_2(a^2b^2c^2))}.
\]
In particular,
\begin{equation}\label{eq:bounds}
 \den(j)\ge\frac{H^{4d}}{256},\qquad
 h(j)\le6d\log H+\log6912.
\end{equation}
Consequently
\begin{equation}\label{eq:exponential}
 \log\den(j)\ge\frac23h(j)
                  -\frac23\log6912-\log256.
\end{equation}
\end{proposition}
\begin{proof}
The Weierstrass invariants are
\[
 c_4=16(a^2+ab+b^2),\qquad
 \Delta=16a^2b^2(a+b)^2,
\]
which give the formula. The positive integers $a,b,c$ are pairwise
coprime. If an odd prime divides $a$, the expression $a^2+ab+b^2$
reduces to $b^2$ and is a unit; the analogous statement holds for $b$.
If it divides $c$, then $a\equiv-b$ and the same expression again
reduces to $b^2$. Thus no odd prime cancels. The expression is odd:
if one of $a,b$ is even, the other is odd, and if both are odd its
three summands are odd. Hence the exact cancellation is the displayed
power of two.

Since $|AB|\ge H$ and $c\ge H^d$, the original denominator is at
least $H^{4d}$, proving the first bound. The numerator is at most
$256(3H^{2d})^3=6912H^{6d}$, and the denominator is at most
$H^{4d}(2H^d)^2=4H^{6d}$. Reduction cannot increase either number,
which proves the second bound. Substitute the resulting lower bound
for $\log H$ into the first bound to obtain~\eqref{eq:exponential}.
\end{proof}

\begin{corollary}\label{cor:finite}
For fixed constants $C>0$ and $k\ge0$, only finitely many primitive
pairs $(A,B)$ satisfy
\[
                    \den(j(E_{A,B,d}))\le C h(j(E_{A,B,d}))^k.
\]
\end{corollary}
\begin{proof}
The first bound in~\eqref{eq:bounds} shows that $h(j)$ tends to
infinity with $H$. The exponential lower bound
\eqref{eq:exponential} eventually exceeds every fixed polynomial
in $h(j)$. Thus $H$ is bounded.
\end{proof}

Pasten's restricted Szpiro theorem~\cite[Theorem 1.2]{PastenSmall}
assumes a bound for $\den(j)$ by a fixed power of
$\log|\operatorname{num}(j)|$. The height-based variant in the attached
survey is weaker than that hypothesis. Corollary~\ref{cor:finite}
therefore excludes both hypotheses on an unbounded family of these
Frey curves. Passing to another Weierstrass model or taking a quadratic
twist leaves $j$ unchanged and does not repair this problem.

\section{Application to the elliptic scalar}

Let $P$ be the fixed nontorsion point of the construction over
$K=\Q(\sqrt2)$, and put
\[
             T(n)=\Tr_{K/\Q}x([n]P),\qquad n\ne0.
\]
The retained trace-height theorem gives constants $A_0>0,J_0\ge0$
such that
\begin{equation}\label{eq:trace}
                     n^2\le A_0h(T(n))+J_0.
\end{equation}
Also $T(n)>6$. At an integer input the scalar pair in the current
construction is $(\beta,T(n))=(n,T(n))$. Write
$n/T(n)=A_n/B_n$ in lowest terms. Then
\[
 \log\max(|A_n|,|B_n|)=h(n/T(n))
 \ge h(T(n))-\log|n|
 \ge\frac{n^2-J_0}{A_0}-\log|n|.
\]
This tends to infinity. Corollary~\ref{cor:finite} consequently
applies already along the required integer-input witnesses. The
restriction is not confined to unwanted rational-input solutions.
The argument applies to every unbounded integer polynomial input
sequence, including both current shifts of squares and shifts of
fourth powers.

\section{The asymmetric-size restrictions}

Theorem 1.3 of Pasten--Sep\'ulveda-Manzo~\cite{PastenSepulveda} applies
to primitive positive triples $a+b=c$ satisfying, for a fixed $\tau>0$,
\begin{equation}\label{eq:asymmetric}
 \min(a,b)\le c\exp\bigl(- (\log c)^\tau\log_2^*c\bigr).
\end{equation}
The earlier restricted estimate recalled in their Theorem 1.2 instead
uses $\min(a,b)\le c^{1-\eta}$ for a fixed $\eta>0$.
The word ``fixed'' is essential for a uniform height conclusion.

\begin{lemma}\label{lem:boundedreal}
Let $p\in\Z[X]$ be nonconstant, tend to positive infinity, and satisfy
$p(0)\ne0$. There are arbitrarily large integers $m$ for which
\[
                    6<T(p(m))\le C_p
\]
with a constant $C_p$ independent of $m$.
\end{lemma}
\begin{proof}
Consider the point $(P,\sigma P)$ in the compact real Lie group
$E(\mathbb R)\times E^\sigma(\mathbb R)$, and the closure $G$ of its
cyclic subgroup. The identity component $G^0$ is a torus, and $G/G^0$
is finite. Choose a positive integer $q$ killing this component group;
the cyclic group generated by $q(P,\sigma P)$ is dense in $G^0$.
For $m=qk$, the polynomial difference $p(qk)-p(0)$ is divisible by $q$.
For each nontrivial character of $G^0$, its value on
$q(P,\sigma P)$ is an irrational rotation, by density and connectedness.
Weyl's polynomial equidistribution theorem therefore shows that
$[p(qk)](P,\sigma P)$ is equidistributed in the coset
$[p(0)](P,\sigma P)+G^0$.

The point $[p(0)]P$ and its conjugate are finite nontorsion points.
Choose a relatively open neighborhood of this tuple whose compact
closure avoids the two identity divisors. The trace of the abscissas
is bounded above on that closure. The neighborhood has positive Haar
measure in the coset and is visited infinitely often. Positivity gives
the lower bound six at all the resulting nonzero multiples.
\end{proof}

\begin{proposition}\label{prop:asymmetric}
Fix a positive even $d$ and a polynomial input sequence as in
Lemma~\ref{lem:boundedreal}. Along infinitely many required integer
witnesses, the primitive triple
\[
             a=A_n^d,\qquad b=B_n^d,\qquad c=A_n^d+B_n^d,
             \qquad A_n/B_n=n/T(n),
\]
fails~\eqref{eq:asymmetric} for every fixed $\tau>0$, eventually
along that subsequence. It also fails
$\min(a,b)\le c^{1-\eta}$ for every fixed $\eta>0$.
\end{proposition}
\begin{proof}
Take the subsequence where $6<T(n)\le C_p$. For sufficiently large
$n$, $A_n>B_n>0$, and
\[
 \log\frac{c}{\min(a,b)}
 =\log\left(1+\left(\frac{n}{T(n)}\right)^d\right)=O(\log n).
\]
By~\eqref{eq:trace} and the preceding lower bound for $h(n/T(n))$,
\[
 \log c\ge d\log A_n
 \ge d\left(\frac{n^2-J_0}{A_0}-\log n\right)\gg n^2.
\]
For every fixed $\tau>0$, the quantity
$(\log c)^\tau\log_2^*c$ grows faster than $\log n$.
This contradicts the logarithmic form of~\eqref{eq:asymmetric}.
Likewise $\eta\log c\gg n^2$ exceeds the displayed logarithmic
ratio, proving the second assertion.
\end{proof}

The proposition rules out a direct uniform invocation of those
restricted estimates on the present scalar triples. It does not
exclude changing the scalar, using a different auxiliary family, or
proving a new restricted inequality with hypotheses adapted to the
actual point and pole equations. It makes no assertion that any of
these triples violates the ordinary $abc$ conjecture.

\section{Exact verification}

The companion script \texttt{check\_frey\_denominator.py} computes the
Weierstrass invariants independently, verifies the rational formula
and its exact two-adic cancellation on a finite exhaustive set, and
checks the integer forms of the two displayed height majorants.
The arguments proving all-input statements and infinite subsequences
are the proofs above; finite enumeration is not substituted for them.

\begin{thebibliography}{9}
\bibitem{PastenSmall}
H.~Pasten, \emph{Szpiro's conjecture when the denominator of the
$j$-invariant is small}, CUBO \textbf{28} (2026), no.~2, 383--389,
\url{https://doi.org/10.56754/0719-0646.2802.383}.
Primary manuscript: \url{https://arxiv.org/abs/2308.07114v2},
Theorem~1.2 and Corollary~3.4.
\bibitem{PastenSepulveda}
H.~Pasten and R.~Sep\'ulveda-Manzo, \emph{On the $abc$ and the $abcd$
conjectures}, Bull.~Braz.~Math.~Soc.~(N.S.) \textbf{56} (2025), article~54,
\url{https://doi.org/10.1007/s00574-025-00476-1}.
Primary manuscript: \url{https://arxiv.org/abs/2406.05083v1},
Theorems~1.2 and~1.3. Only their precise size hypotheses are used here.
\bibitem{Weyl}
H.~Weyl, \emph{"Uber die Gleichverteilung von Zahlen mod.~Eins},
Math.~Ann.~\textbf{77} (1916), 313--352,
\url{https://doi.org/10.1007/BF01475864}.
\bibitem{Trace}
Project repository, \emph{The height of an aggregate elliptic trace},
\texttt{aggregate\_trace\_height\_20261108\_v9.tex}, one-point case.
\end{thebibliography}
\end{document}
